% ---------------------------------------------------------------------------
% SAR-YOLO — CVPR-style manuscript skeleton
%
% IMPORTANT: every table and figure input in this file is GENERATED by
%   python -m saryolo assets
% and contains TBD for any cell that has not been measured. Do not hand-fill a
% TBD cell: run the experiment, or state in the text that it was not run.
%
% The claim-checking discipline for this paper (see ../README.md):
%   * the ablation table adds one component per row;
%   * `att_saa_static` and `fus_static` are the runs that justify the word
%     "adaptive" -- if they match the proposed blocks, the claim must be dropped;
%   * every reported number must be traceable to a ledger run id.
% ---------------------------------------------------------------------------
\documentclass[10pt,twocolumn,letterpaper]{article}

\usepackage{cvpr}              % replace with the official CVPR style file
\usepackage{graphicx}
\usepackage{amsmath,amssymb}
\usepackage{booktabs}
\usepackage{multirow}
\usepackage{xcolor}
\usepackage[pagebackref,breaklinks,colorlinks]{hyperref}

\newcommand{\tbd}{\textcolor{red}{TBD}}
\newcommand{\ours}{\textsc{SAR-YOLO}}

\title{\ours: Speckle-Aware Feature Representation for\\
Robust Object Detection in SAR Imagery}

\author{Arghya Bose\\
\small Affiliation\\
\small \texttt{officialarghya29@gmail.com}
}

\begin{document}
\maketitle

% ===========================================================================
\begin{abstract}
Synthetic aperture radar (SAR) provides all-weather, day-night imaging, but its
imagery is dominated by multiplicative speckle, background clutter and low
contrast, and the targets of interest are frequently small and span a wide scale
range. Standard detectors, developed for natural images, lose precisely the
target evidence that SAR preserves.
%
We introduce \ours{}, a SAR-specific object detector built on a YOLO baseline
through four individually ablatable components: a \emph{SAR Feature Enhancement}
block that re-expresses features through local contrast normalisation rather than
destructive denoising; a \emph{Speckle-Aware Feature Module} that estimates a
per-channel speckle likelihood and suppresses it multiplicatively, respecting
the signal-dependent nature of speckle; a \emph{SAR-Adaptive Attention} block
that fuses channel, spatial and local-contrast evidence with input-dependent
weights; and an \emph{Adaptive Multi-Scale Fusion} block that reweights
concatenated neck branches per sample. A high-resolution P2 detection head and a
SAR-aware auxiliary objective complete the model.
%
Every component is an exact identity at initialisation, so \ours{} starts
numerically identical to its baseline and each reported gain is learned
behaviour rather than added initial capacity. We report accuracy, scale-wise AP,
robustness under controlled corruption, efficiency, cross-dataset generalisation
and multi-seed stability.
%
% NOTE: do not write a quantitative claim here until the corresponding table is
% free of TBD cells.
An initial experimental campaign on \tbd{dataset} yields \tbd{mAP50} against
\tbd{baseline} for the YOLO baseline.
\end{abstract}

% ===========================================================================
\section{Introduction}
\label{sec:intro}

SAR imaging is indispensable when optical sensors fail: it penetrates cloud and
operates without illumination. Its statistics, however, are hostile to detectors
designed for photographs. Speckle is multiplicative and signal-dependent;
background clutter produces bright, target-like responses; contrast is low; and
targets are often only a few pixels across while simultaneously spanning a wide
scale range within a single scene.

Existing detectors inherit design decisions from natural images. Pooling and
fixed-weight feature fusion assume that a strong response is informative, which
is false under speckle. Attention blocks such as SE, ECA and CBAM fuse their
branches with fixed weights that cannot adapt to the imaging regime actually
present. Plain concatenation in a PAN-FPN neck gives every scale the same
influence, so a small, speckle-obscured target can be swamped by a large
high-energy neighbour.

This paper asks a narrow question: \emph{can a small number of SAR-specific,
evidence-driven modifications to a standard detector improve detection of weak
and small targets without materially increasing cost?} We answer it with a
controlled study rather than a component stack:
Section~\ref{sec:method} derives each module from a measured SAR property,
Section~\ref{sec:experiments} evaluates them one at a time, and
Section~\ref{sec:ablation} isolates each contribution.

\paragraph{Contributions.}
\begin{itemize}
  \item A SAR-aware feature representation that uses local first- and
        second-order statistics to suppress clutter while preserving target
        structure, formulated as a gated residual so it cannot destroy the
        feature it augments (Sec.~\ref{sec:sfe},~\ref{sec:sfm}).
  \item A lightweight attention block whose branch fusion weights are predicted
        per sample, together with the static-weight control experiment that
        isolates the value of adaptivity (Sec.~\ref{sec:saa}).
  \item An adaptive multi-scale fusion operator that reweights concatenated neck
        branches per sample, with a static-weight control (Sec.~\ref{sec:amf}).
  \item A comprehensive evaluation protocol covering ablation, scale-wise AP,
        controlled robustness, efficiency, cross-dataset generalisation and
        multi-seed stability (Sec.~\ref{sec:experiments}).
\end{itemize}

\paragraph{A note on honesty.} Every component is initialised as an exact
identity, and each is compared against the standard component it replaces in the
same architectural slot. Where an experiment does not support a claim, we say so.

% ===========================================================================
\section{Related Work}
\label{sec:related}

\paragraph{SAR object detection.}
Early SAR detectors relied on hand-crafted features such as CFAR and
texture statistics. Modern work applies CNN and transformer detectors to SAR
benchmarks \cite{ssdd,hrsid,sardet100k,sarship}. The dominant approach transfers
a generic detector and modifies the data or the training schedule; comparatively
few works modify the feature representation to respect SAR statistics, and fewer
still isolate the contribution of an individual modification with a control.

\paragraph{YOLO-based detection.}
\cite{yolo,ultralytics} provide the real-time single-stage baseline used here.
We adopt the architecture unmodified as the reference point and measure every
change against it under an identical training pipeline.

\paragraph{Small-object detection.}
High-resolution feature levels (P2) and feature-pyramid rebalancing are standard
remedies. Because adding a P2 level is the most expensive change in our model, we
gate it on a measurement: the COCO size distribution of the dataset.

\paragraph{Speckle and SAR representation.}
Classical despeckling (Lee, Frost, Kuan) reduces speckle but attenuates the
high-frequency evidence that carries target signal, and is applied before the
network rather than learned within it. We instead learn a speckle-likelihood map
and suppress it multiplicatively, in a formulation that reduces to the identity
at initialisation.

\paragraph{Attention.}
SE \cite{se}, ECA \cite{eca} and CBAM \cite{cbam} are the baselines our attention
slot is compared against. Our contribution is not the existence of attention but
the input-dependent fusion of three evidence types, plus the control that tests
whether adaptivity itself matters.

\paragraph{Multi-scale fusion.}
FPN/PAN \cite{fpn} and BiFPN-style weighted fusion motivate our AMF block. The
distinction we test is adaptivity: static learned weights versus per-sample
weights.

% ===========================================================================
\section{Method}
\label{sec:method}

\subsection{Design constraints}
Every module is written as $F' = F + \alpha\,(g(F) - F)$ with
$\alpha = \tanh(\cdot)$ initialised to zero, so the module is an exact identity
before training and a freshly built \ours{} predicts identically to its baseline.
Each module is also channel-preserving and single-input, which lets it be
inserted into a standard YOLO graph without special-casing the model parser.

\subsection{Shared SAR primitive: local statistics}
For feature map $F$ and a $k\times k$ average pool $A$,
\begin{equation}
\mu = A(F),\quad
\sigma^2 = \max\!\big(A(F^2) - \mu^2,\,0\big),\quad
S = \frac{F-\mu}{\sqrt{\sigma^2+\epsilon}}.
\end{equation}
$S$ is a differentiable local contrast normalisation; $\sigma$ measures local
speckle strength. Every component below consumes these statistics, because a
$3\times3$ convolution over raw intensities cannot distinguish uniform clutter
from a structured target.

\subsection{SAR Feature Enhancement (SFE)}
\label{sec:sfe}
\begin{equation}
Z = \phi\big([\,F \,;\, S \,;\, \sigma\,]\big),\qquad
F' = F \odot \gamma + \beta + \alpha\, Z .
\end{equation}
Enhancement is additive; at initialisation $\gamma=1,\beta=0,\alpha=0$ it is the
identity. Alternatives compared in the same slot: identity, log compression,
pure local standardisation, and a differentiable CLAHE analogue.

\subsection{Speckle-Aware Feature Module (SFM)}
\label{sec:sfm}
\begin{equation}
N = \sigma\!\big(\psi([\,F \,;\, \mathrm{hp} \,;\, \sigma^2\,])\big),\quad
T = \mathrm{dwconv}(F),\quad
F' = F + \alpha\,(T - \beta\, N \odot F),
\end{equation}
with $\mathrm{hp} = F-\mu$ and $\beta = \mathrm{sigmoid}(\cdot)$. The suppression
term is multiplicative in $F$ because speckle intensity scales with reflectivity;
a constant offset would suppress weak and strong regions by the same absolute
amount, which is the wrong behaviour. $\beta$ is constrained to $(0,1)$ so the
module cannot learn to amplify speckle.

\subsection{SAR-Adaptive Attention (SAA)}
\label{sec:saa}
Let $d(F) = [\overline{F}_c ; \max_c F ; \mathrm{std}_c F]$ be a channel
descriptor, and let
\begin{align}
a &= \sigma(\mathrm{MLP}_c(d(F))),\\
s &= \sigma(\mathrm{conv}([\overline{F}_c ; \max_c F ; \sigma^2 ; \sigma])),\\
\ell &= \sigma(\psi([\,|\overline{S}|_c ; |\overline{\mathrm{hp}}|_c\,])),\\
w &= \mathrm{softmax}(\mathrm{MLP}_w(d(F))),\qquad M = w_0 a + w_1 s + w_2 \ell,\\
F' &= F + \alpha\,(F \odot M - F).
\end{align}
Two properties separate this from CBAM: $\mathrm{std}_c F$ enters the descriptor,
since speckle strength is a second-order statistic; and $w$ is predicted
\emph{per sample}, so the network can change emphasis with the imaging regime.
The control experiment (\texttt{att\_saa\_static}) replaces $w$ with learned
constants shared across the batch; if it matches, the adaptivity claim does not
hold.

\subsection{Adaptive Multi-Scale Fusion (AMF)}
\label{sec:amf}
Inserted after each neck concatenation, with the concatenated channels split into
groups $F_1,\dots,F_G$ (one per fused scale):
\begin{equation}
w = \mathrm{softmax}\big(\mathrm{MLP}([d(F_1);\dots;d(F_G)])\big),\quad
\tilde{F} = [\,w_1F_1;\dots;w_GF_G\,],
\end{equation}
\begin{equation}
M = \sigma\big(\mathrm{dwconv}(\psi(\tilde F))\big),\qquad
F' = F + \alpha\,(\tilde F \odot M - F).
\end{equation}
Group sizes are resolved from the parser's channel list, so the split is correct
at any width multiplier; the module validates that the groups sum to its input
width, so a stale configuration fails loudly. A plain additive fusion baseline is
not defined here because PAN-FPN branches have unequal widths, so the additive
baseline 1$\times$1-projects each branch to a shared width before summing, which
is how residual-Add fusion is implemented when widths differ.

\subsection{Small-object detection head}
A stride-4 P2 level is added to the neck. This is the most expensive change in
the model, so it is included only if the dataset's object-size distribution is
dominated by COCO-small instances; the measurement is reported in
Sec.~\ref{sec:dataset}.

\subsection{SAR-aware objective}
\begin{equation}
\mathcal{L} = \mathcal{L}_{\mathrm{box}} + \mathcal{L}_{\mathrm{cls}}
+ \mathcal{L}_{\mathrm{dfl}}
+ \lambda_s \mathcal{L}_{\mathrm{sep}}
+ \lambda_o \mathcal{L}_{\mathrm{small}}
+ \lambda_r \mathcal{L}_{\mathrm{smooth}} .
\end{equation}
$\mathcal{L}_{\mathrm{sep}}$ is a margin ranking between mean foreground
confidence and the hardest background anchors,
$\mathrm{relu}(m - (\overline{c}_{fg} - \overline{c}_{bg}^{\text{top}}))$, so the
model is rewarded for separating targets from clutter rather than for being
uniformly unconfident. $\mathcal{L}_{\mathrm{small}}$ adds a confidence
objective restricted to anchors assigned to small ground-truth boxes, and
$\mathcal{L}_{\mathrm{smooth}}$ penalises total variation of the
background-masked score map to suppress speckle-induced scattered activations.
With all weights at zero the objective is numerically identical to the baseline
loss, which is what makes the loss ablation row interpretable.

% ===========================================================================
\section{Experiments}
\label{sec:experiments}

\subsection{Dataset}
\label{sec:dataset}
\tbd{Insert dataset statistics from \texttt{dataset\_statistics/}}: image count,
class distribution, objects per image, object-size distribution (COCO
small/medium/large), aspect ratios, object density heat map, and the SAR
difficulty proxies (mean local contrast and target/background intensity ratio).
The size distribution must be reported here because it is what justifies or
rejects the P2 head.

\subsection{Experimental setup}
All models are trained with an identical pipeline, schedule, augmentation and
input resolution; results are compared at the same input size. Seeds, versions
and hardware are recorded per run in \texttt{results/experiments.jsonl}.

\subsection{Baseline comparison}
\begin{table}[t]
\centering
% WARNING: contains unmeasured (TBD) cells -- not submission ready.
\caption{Comparison with YOLO baselines and the proposed SAR-YOLO. Cells are TBD until the corresponding experiment has been run.}
\label{tab:baseline_comparison}
\begin{tabular}{lrrrrrrr}
\toprule
Model & mAP50 & mAP50:95 & Precision & Recall & Params (M) & GFLOPs & FPS \\
\midrule
YOLO11 baseline & TBD & TBD & TBD & TBD & TBD & TBD & TBD \\
SAR-YOLO (ours) & TBD & TBD & TBD & TBD & TBD & TBD & TBD \\
\bottomrule
\end{tabular}
\end{table}


\subsection{Main ablation}
\label{sec:ablation}
\begin{table}[t]
\centering
% WARNING: contains unmeasured (TBD) cells -- not submission ready.
\caption{Main ablation. Each row adds exactly one component, so every delta is attributable to that component alone.}
\label{tab:main_ablation}
\begin{tabular}{lrrrrrrrrrr}
\toprule
Model & SFE & Speckle & Attention & AMF & P2 head & SAR loss & mAP50 & mAP50:95 & Params (M) & FPS \\
\midrule
YOLO baseline &  &  &  &  &  &  & TBD & TBD & TBD & TBD \\
+ SFE & \checkmark &  &  &  &  &  & TBD & TBD & TBD & TBD \\
+ Speckle & \checkmark & \checkmark &  &  &  &  & TBD & TBD & TBD & TBD \\
+ Attention & \checkmark & \checkmark & \checkmark &  &  &  & TBD & TBD & TBD & TBD \\
+ AMF & \checkmark & \checkmark & \checkmark & \checkmark &  &  & TBD & TBD & TBD & TBD \\
+ Small head & \checkmark & \checkmark & \checkmark & \checkmark & \checkmark &  & TBD & TBD & TBD & TBD \\
Full & \checkmark & \checkmark & \checkmark & \checkmark & \checkmark & \checkmark & TBD & TBD & TBD & TBD \\
\bottomrule
\end{tabular}
\end{table}


\subsection{Module-level ablation}
\begin{table}[t]
\centering
% WARNING: contains unmeasured (TBD) cells -- not submission ready.
\caption{Module-level comparison. Each proposed block is compared against the standard component it replaces, in the same architectural slot, so the comparison isolates the mechanism rather than the added parameters.}
\label{tab:module_ablation}
\begin{tabular}{lrrrr}
\toprule
Slot & Variant & mAP50 & mAP50:95 & Params (M) \\
\midrule
Attention & att_none & TBD & TBD & TBD \\
Attention & att_se & TBD & TBD & TBD \\
Attention & att_eca & TBD & TBD & TBD \\
Attention & att_cbam & TBD & TBD & TBD \\
Attention & att_saa_static & TBD & TBD & TBD \\
Attention & attention & TBD & TBD & TBD \\
Fusion & fus_concat & TBD & TBD & TBD \\
Fusion & fus_add & TBD & TBD & TBD \\
Fusion & fus_static & TBD & TBD & TBD \\
Fusion & amf & TBD & TBD & TBD \\
Preprocessing & pre_identity & TBD & TBD & TBD \\
Preprocessing & pre_log & TBD & TBD & TBD \\
Preprocessing & pre_clahe & TBD & TBD & TBD \\
Preprocessing & pre_standardize & TBD & TBD & TBD \\
Preprocessing & sfe & TBD & TBD & TBD \\
Speckle & spk_none & TBD & TBD & TBD \\
Speckle & spk_lee & TBD & TBD & TBD \\
Speckle & spk_denoise & TBD & TBD & TBD \\
Speckle & speckle & TBD & TBD & TBD \\
\bottomrule
\end{tabular}
\end{table}


\subsection{Scale-wise analysis}
\begin{table}[t]
\centering
% WARNING: contains unmeasured (TBD) cells -- not submission ready.
\caption{Performance by object scale. Scale-wise AP is reported separately because a single mAP can hide exactly the small-object effect this work claims. 'n GT' is the number of ground-truth boxes in each range; a range with none is unmeasurable, not zero.}
\label{tab:scale_analysis}
\begin{tabular}{lrrrrrr}
\toprule
Model & AP_small & n GT small & AP_medium & n GT medium & AP_large & n GT large \\
\midrule
YOLO baseline & TBD & TBD & TBD & TBD & TBD & TBD \\
SAR-YOLO (ours) & TBD & TBD & TBD & TBD & TBD & TBD \\
\bottomrule
\end{tabular}
\end{table}

A range with no ground truth is reported as not measurable rather than as zero,
because a zero would be read as a model failure.

\subsection{Robustness}
\begin{table}[t]
\centering
% WARNING: contains unmeasured (TBD) cells -- not submission ready.
\caption{Robustness under controlled degradations. Both models face byte-identical corrupted inputs, and only the images are degraded (ground truth is untouched).}
\label{tab:robustness}
\begin{tabular}{lrrrr}
\toprule
Corruption & Severity & YOLO mAP50 & SAR-YOLO mAP50 & Drop (ours) \\
\midrule
TBD & TBD & TBD & TBD & TBD \\
\bottomrule
\end{tabular}
\end{table}


\subsection{Efficiency}
\begin{table}[t]
\centering
% WARNING: contains unmeasured (TBD) cells -- not submission ready.
\caption{Efficiency. Latency is measured after warm-up with CUDA synchronisation; FLOPs are reported at the recorded input size, since FLOPs are meaningless without it.}
\label{tab:efficiency}
\begin{tabular}{lrrrrr}
\toprule
Model & Params (M) & GFLOPs & FPS & Latency (ms) & Size (MB) \\
\midrule
YOLO baseline & TBD & TBD & TBD & TBD & TBD \\
SAR-YOLO (ours) & TBD & TBD & TBD & TBD & TBD \\
\bottomrule
\end{tabular}
\end{table}


\subsection{Multi-seed stability}
\begin{table}[t]
\centering
% WARNING: contains unmeasured (TBD) cells -- not submission ready.
\caption{Multi-seed validation. A single-seed difference of a few tenths of a point is not evidence, so the headline result is reported as mean +/- std over seeds.}
\label{tab:multi_seed}
\begin{tabular}{lrrrr}
\toprule
Metric & Mean & Std & n seeds & Per-seed values \\
\midrule
mAP50:95 & TBD & TBD & 0 & TBD \\
\bottomrule
\end{tabular}
\end{table}


\subsection{Qualitative results and failure analysis}
Qualitative panels compare ground truth, baseline and \ours{} on dense, sparse
and small-object scenes; Grad-CAM and feature maps support the claim that the
modules concentrate activations on target structure. Failure cases are
categorised explicitly (false positive, false negative, localisation error,
classification error, small-object miss, clutter confusion) and reported with
counts rather than hidden.

% ===========================================================================
\section{Conclusion}
\label{sec:conclusion}
We presented \ours{}, a SAR-specific detector built by adding four individually
ablatable components to a YOLO baseline. The design rule throughout was that
each component addresses a measured SAR property, is an exact identity at
initialisation, and is compared against the standard component it replaces in the
same slot. \tbd{Insert the supported conclusions once the tables are complete;
state explicitly which claims the control experiments failed to support.}

\paragraph{Limitations.}
\tbd{State any component whose control experiment did not show a significant
benefit, the datasets on which the method was not evaluated, and the
scene-leakage caveat for chip-based SAR datasets.}

{\small
\bibliographystyle{ieee_fullname}
\bibliography{refs}
}

\end{document}
